\section{Derivative calculation and component bookkeeping}

For cyclic spatial structure constants and aligned torsion, the nonzero connection coefficients, together with antisymmetry in the first two indices, are

\begin{equation}\tag{B-01}\label{eq:B-01}
\Gamma_{ijk}=\left(\frac{n_i+n_j-n_k}{2}-h\right)\varepsilon_{ijk},
\qquad \Gamma_{i3j}=-\frac V3\delta_{ij},\qquad
\Gamma_{3ij}=\frac V3\delta_{ij}.
\end{equation}

The full noncoordinate-frame curvature is

\begin{equation}\tag{B-02}\label{eq:B-02}
\begin{aligned}
R_{abcd}={}&D_c\Gamma_{abd}-D_d\Gamma_{abc}
+\Gamma_{aec}\Gamma_{ebd}-\Gamma_{aed}\Gamma_{ebc}
+\Gamma_{abe}C^e{}_{cd}.
\end{aligned}
\end{equation}

The first two terms retain the derivatives of the connection, including contortion. For circle profiles, projection of (B-02) gives the LC Weyl tensor and the following additional entries:

\begin{equation}\tag{B-03}\label{eq:B-03}
\begin{aligned}
a_+&=\tfrac12\left(\dot h+\tfrac23hV\right),\qquad
a_-=\tfrac12\left(\dot h-\tfrac23hV\right),\\
R^{(2)}_{ijk3}&=a_+\varepsilon_{ijk},\qquad
R^{(2)}_{k3ij}=-a_+\varepsilon_{ijk},\\
R^{(3)}_{abcd}&=a_-\varepsilon_{abcd}.
\end{aligned}
\end{equation}

All other entries follow from antisymmetry within pairs, or vanish. Summing the twenty-four nonzero signed entries in each component gives $24a_+^2$ and $24a_-^2$, proving (D-01). Contracting the unprojected curvature gives the scalar term in (D-01); contracting its four-form part gives $12a_-=6\dot h-4hV$ in (D-06). A separate exterior-algebra construction verifies the first Cartan equation, both differential Bianchi identities, and the Nieh--Yan identity.

The Pontryagin density in (D-06) follows either from the exterior product of the curvature two-forms or from the connection Chern--Simons form. With the matrix-trace convention, $\mathrm{Tr}(R\wedge R)=-\sum_{a,b}R_{ab}\wedge R_{ab}$, so $\mathrm{d}\mathrm{CS}(\omega)=-\mathcal P\,\mathrm{vol}$. This also fixes the sign of the displayed primitive.

To expose the cancellation between the two sectors, introduce separate symbols $b_2,b_3$ only as bookkeeping weights. Their bulk contribution would read

\begin{equation}\tag{B-04}\label{eq:B-04}
-b_2I_2-b_3I_3
=-6(b_2+b_3)\dot h^2-8(b_2-b_3)hV\dot h
-\frac83(b_2+b_3)h^2V^2.
\end{equation}

The action compared in this paper has $b_2=b_3=\alpha$, not independent physical couplings. Moreover, $I_2-I_3=16hV\dot h$ need not integrate to zero when $V$ varies. In physical circle coordinate $z$,

\begin{equation}\tag{B-05}\label{eq:B-05}
\int_0^L(I_2-I_3)\,\mathrm{d}z
=-8\int_0^L\dot V h^2\,\mathrm{d}z.
\end{equation}

The cancellation in (D-02) is therefore a consequence of equal weights, not a general claim that the two component functionals are identical. The reduction to the scalar action at $\alpha=0$ and the static Paper II limit at $\dot h=\dot V=0$ are separate checks.
